% answers_short.tex —— 1993 年·新科目组·理科 · 简短答案速查
% 由 tools/make_short_answers.py 自 answers.tex 的 \gksolhead/\gkanswer 提取，
% 并与 1993/1993_数学真题_新科目组理_答案要点.md 逐题比对后生成，勿手改（改源 MD 后重跑）。

\begin{gknotice}
\textbf{使用说明}：本节只给最终答案便于快速核对，完整推导见解析版。答案为编写者独立推导，\textbf{非官方标准答案}。
\end{gknotice}

\gkpart{选择题}
\begin{center}\begin{tabular}{cccccccccccccccccc}
\toprule
题号 & 1 & 2 & 3 & 4 & 5 & 6 & 7 & 8 & 9 & 10 & 11 & 12 & 13 & 14 & 15 & 16 & 17 \\
\midrule
答案 & A & C & B & B & A & D & C & A & B & D & C & A & D & D & A & C & B \\
\bottomrule
\end{tabular}\end{center}

\gkpart{填空题}
\begin{itemize}[leftmargin=2.6em]
  \item[\textbf{18.}] \(\dfrac{\sqrt3+2\sqrt2}{6}\)
  \item[\textbf{19.}] \(|k|>\dfrac13\)
  \item[\textbf{20.}] \(100\)
  \item[\textbf{21.}] \(\{1\}\)（按 \(f^{-1}(0)\) 为原像/零点集理解）
  \item[\textbf{22.}] \(1760\) 元
  \item[\textbf{23.}] \(30^\circ\)
\end{itemize}

\gkpart{解答题}
\begin{itemize}[leftmargin=2.6em]
  \item[\textbf{24.}] (1) 定义域 \((-1,1)\)；(2) 奇函数（证明见解析）；(3) \(a>1\) 时 \(x\in(0,1)\)；\(0<a<1\) 时 \(x\in(-1,0)\)
  \item[\textbf{25.}] \(S_n=1-\dfrac{1}{(2n+1)^2}\)
  \item[\textbf{26.}] 两个结论均成立（两平面互相垂直、两直线互相垂直）
  \item[\textbf{27.}] \(\dfrac{x^2}{15/4}+\dfrac{y^2}{3}=1\)
  \item[\textbf{28.}] 原题无唯一解集：主值 \((-\pi,\pi]\) 下四个候选均满足，\([0,2\pi)\) 下仅 \(\theta=\dfrac{\pi}{12},\dfrac{7\pi}{12}\)
\end{itemize}
